
% ============================================================
% Packages
% ============================================================

\usepackage{microtype}
\usepackage{graphicx}
\usepackage{booktabs}
\usepackage{longtable}
\usepackage{xcolor}
\usepackage{tikz}
\usepackage{titlesec}
\usepackage{tocloft}
\usepackage{babel}
\babelprovide[
  main,
  import,
  transforms = punctuation.space
]{french}



% ============================================================
% Table des matières
% ============================================================

% Parties en gras
\renewcommand{\cftpartfont}{\bfseries}
\renewcommand{\cftpartpagefont}{\bfseries}

% Chapitres en romain
\renewcommand{\cftchapfont}{\normalfont}
\renewcommand{\cftchappagefont}{\normalfont}

% Espacements
\setlength{\cftbeforechapskip}{2pt}
\setlength{\cftbeforepartskip}{12pt}


% ============================================================
% Polices
% ============================================================

\usepackage{fontspec}

% Texte courant
%\setmainfont{TeX Gyre Heros}[
\setmainfont{TeX Gyre Adventor}[
  Ligatures=TeX
]


% Police des titres
\newfontfamily\titlefont{TeX Gyre Pagella}

% ============================================================
% Titres de chapitre
% ============================================================

% Parties
\titleformat{\part}[display]
  {\titlefont\Huge\bfseries}
  {\centering\partname~\thepart}
  {12pt}
  {}

% Chapitres
\titleformat{\chapter}[display]
  {\titlefont}
  {\large\chaptername~\thechapter}
  {8pt}
  {\Huge\bfseries}

% Sections
\titleformat{\section}
  {\titlefont\Large\bfseries}
  {\thesection.}
  {0.5em}
  {}

% espace avant / après le titre
\titlespacing*{\chapter}
  {0pt}
  {0pt}
  {6pt}


\renewcommand{\thesection}{\arabic{section}}

\newcommand{\chapterrunninghead}[1]{%
  \markboth{\chaptername~\thechapter}{\chaptername~\thechapter}%
}


% ============================================================
% Auteur d'un chapitre
% ============================================================

\newcommand{\chapterauthor}[1]{%
  {\large\itshape #1\par}%
  \vspace{24pt}%
}


% ============================================================
% Page de couverture
% ============================================================

\newcommand{\covertitle}[1]{%
  {\fontsize{24}{24}\selectfont\bfseries #1\par}%
}

\newcommand{\coversubtitle}[1]{%
  {\LARGE #1\par}%
}

\newcommand{\coverdate}[1]{%
  {\LARGE #1\par}%
}

% ============================================================
% Numérotation chapitres
% ============================================================

\setlength{\cftchapnumwidth}{3em}

% Numérotation des chapitres par partie : I.1, I.2, II.1...
\makeatletter
\@addtoreset{chapter}{part}
\makeatother

\renewcommand{\thechapter}{\thepart.\arabic{chapter}}

% Identifiant interne unique pour hyperref
\providecommand{\theHchapter}{}
\renewcommand{\theHchapter}{\arabic{part}.\arabic{chapter}}


% ============================================================
% Headers et pieds de page
% ============================================================

\usepackage{fancyhdr}

\pagestyle{fancy}
\fancyhf{}

% Pages paires
\fancyhead[LE]{Chapitre \thechapter}

% Pages impaires
\fancyhead[RO]{Partie \thepart}

% Numéros de page
\fancyfoot[LE,RO]{\thepage}

\setlength{\headheight}{13.6pt}
